\PageTitle{04 / Benchmark detail}{Per-tool latency and returned rows}
Every tool has 20 successful samples. Figure~\ref{fig:tools} compares sample medians and p95 values. Tool names are shortened on the axis; Table~\ref{tab:tools} gives the exact names.

\begin{figure}[h]\centering\begin{tikzpicture}
\begin{axis}[ybar,bar width=9pt,ymin=0,ymax=750,ylabel={Client-observed latency (ms)},xtick={1,2,3,4,5},xticklabels={Items,Sales orders,Purchase orders,Assets,Search Item},x tick label style={font=\sffamily\scriptsize},ymajorgrids=true,legend style={at={(.5,1.04)},anchor=south,legend columns=2},height=70mm,enlarge x limits=.12]
\addplot[fill=ConversationBlue,draw=ConversationBlue] coordinates {(1,337) (2,351) (3,352) (4,370) (5,390)};\addlegendentry{Median}
\addplot[fill=white,draw=TaskAmber,line width=1pt] coordinates {(1,628) (2,643) (3,656) (4,662) (5,634)};\addlegendentry{p95}
\end{axis}
\end{tikzpicture}
\caption{Per-tool client latency. Filled bars are medians; outlined bars are p95. Bars share a zero baseline.}\label{fig:tools}\end{figure}
\Source{S2; nearest rank, 20 calls per tool}

\begin{table}[h]\caption{Per-tool results with exact operation names}\label{tab:tools}\EvidenceTableStart
\begin{tabularx}{\linewidth}{@{}Xrrrr@{}}\toprule
Tool & Calls & Rows & Median & p95\\
 & & / response & ms & ms\\\midrule
\texttt{list\_items} & 20/20 & 56 & 337 & 628\\
\texttt{list\_sales\_orders} & 20/20 & 5 & 351 & 643\\
\texttt{list\_purchase\_orders} & 20/20 & 10 & 352 & 656\\
\texttt{list\_assets} & 20/20 & 3 & 370 & 662\\
\texttt{search\_records(Item)} & 20/20 & 56 & 390 & 634\\\bottomrule
\end{tabularx}\end{table}

\subsection*{What this comparison does and does not show}
The observed medians range from 337 to 390 ms and p95 values from 628 to 662 ms. Collections differ in size and the calls share the same burst run. These summaries do not support a causal claim that one tool implementation is faster or a statistically established ranking of tool performance.

The page cap was 100, but the largest collection returned 56 rows. Only \texttt{name} was requested. This workload did not stress large-result serialization, broad record projection, pagination at the cap, writes, or mixed business transactions.

\KeyPoint{Use the correct unit}{The rows column is the number of records per response, not the number of requests or an independent correctness score. Row counts and success flags do not validate the business meaning of every returned record.}
